\documentclass[exercice]{thedocument}%
\usepackage{texmaths-tikz}
\usepackage{tabularray}
\startdatakeys
\source{originale}
\cycle{4}
\competencesDuSocle{RE}
\connaissancesRequises{B4b*}
\competencesTravaillees{B4b*}
\enddatakeys
\begin{document}%
\!func={xmin=-4,xmax=4,ymin=-2,ymax=2,gridstep=0.25,func=f=x^2-1.5;};%
\!scale=1;%
\edef\func{\*func;}%
\noindent
\begin{minipage}{.5\linewidth}
Voici, ci-contre, la représentation graphique d'une fonction $f$.
Par lecture graphique~:
\begin{enumerate}
    \item Quelle est l'image de \<n1; par la fonction $f$~? Faire une phrase réponse.
    \item Traduire la phrase réponse précédente en notation mathématique.
    \item Donner un antécédent de \<n2; par la fonction $f$~? Faire une phrase réponse.
    \item Traduire la phrase réponse précédente en notation mathématique.
    \item Compléter le tableau suivant :\par\noindent
    \begin{tblr}{colspec=*{6}{Q[c,m,wd=30pt]},column{1}={font=\bfseries,bg=\maintblrbgcolor},hlines,vlines}
        $x$&\<n3;&\<n4;&&\<n6;\\
        $f(x)$&&&\<n5;&&\<n7;\\
    \end{tblr}
\end{enumerate}
\end{minipage}\hfill
\begin{minipage}{.45\linewidth}
    \begin{tikzpicture}[scale=\*scale;]
    \expandafter\simpleFunc\func
\end{tikzpicture}
\end{minipage}\par
\end{document}
